% ====================================================================
\chapter{Sets}\label{ch:sets}
% ====================================================================

\begin{chapterintro}
A set is a collection of objects. Sets are the containers of mathematics: the solutions of an inequality
form a set, the options open to a consumer form a set, the natural numbers form a set. This chapter
introduces the notation for sets, the relations between them (membership, subset, equality) and the
operations that build new sets from old ones (intersection, union and difference).
\end{chapterintro}

\begin{prerequisites}
The language of \cref{ch:logic}: propositions, implication and equivalence (\cref{sec:implication,sec:equivalence}),
and the meaning of ``for all'' and ``there exists'' (\cref{sec:quantifiers}).
\end{prerequisites}

\section{Sets and elements}\label{sec:sets-basic}

\begin{definition}[Set, element][def:set]
A \term{set} is a collection of objects. If the set $A$ consists of the objects $x,y,z,\dots$, we write
$A=\{x,y,z,\dots\}$. If an object $x$ is in the set $A$, we write $x\in A$ and call $x$ an
\term{element} of $A$. If not, we write $x\notin A$. \src{L1, slide 10}
\end{definition}

\begin{sourcenote}[``What is a collection?'']
The slide defines a set as ``a collection (what is a collection?) of objects''. The parenthesis is
deliberate: ``collection'' is no more precise than ``set'', so the definition is not really a definition.
In this module ``set'' is treated as a basic, undefined notion, as it is in most of mathematics, and you
will not be asked to define it from anything simpler. What matters is how sets are \emph{used}: an object
is either an element of a given set or it is not.
\end{sourcenote}

Two features of sets follow from the fact that a set is determined only by \emph{which} objects belong
to it.

\begin{note}[Order and repetition do not matter]
$\{1,2,3\}$, $\{3,1,2\}$ and $\{1,1,2,3,3\}$ are the same set: each contains exactly the elements $1$, $2$
and $3$. Listing an element twice does not put it in the set twice, and the order of the list has no
meaning. (This follows from the definition of equality in \cref{sec:set-equality}.)
\end{note}

\begin{example}[Membership]
Let $A=\{2,4,6,8\}$ and $B=\{\text{apples},\text{oranges}\}$. Then $4\in A$, $5\notin A$,
$\text{apples}\in B$ and $\text{bananas}\notin B$. Elements of a set can be any objects at all, including
other sets: the set $C=\{\{1,2\},\{3\}\}$ has two elements, the set $\{1,2\}$ and the set $\{3\}$. Note
that $1\notin C$: the number $1$ is an element of an element of $C$, not of $C$ itself.
\end{example}

\begin{definition}[Empty set][def:empty]
A set that contains no objects is called the \term{empty set}. It is denoted by the symbol $\emptyset$;
sometimes it is written $\{\}$. \src{L1, slide 10}
\end{definition}

The empty set is the set analogue of the number zero. It appears naturally as the answer to questions
whose answer is ``nothing'': the set of real numbers $x$ with $x^2<0$ is $\emptyset$, and so is the set of
solutions of the inequality $0>1$.

\begin{warning}[$\emptyset$ is not the same as $\{\emptyset\}$]
$\emptyset$ has no elements. $\{\emptyset\}$ is a set with exactly one element, namely the empty set. A box
containing an empty box is not an empty box.
\end{warning}

\subsection{Not everything is a set}
\begin{note}[Russell's paradox]
The slides remark that not everything can be described as a set: for example, there is no ``set of all
sets'' (nor certain similar collections), because assuming that there is leads to a contradiction known
as \emph{Russell's paradox}. \src{L1, slide 10}
\end{note}

\begin{external}[What goes wrong]
Suppose any description at all could define a set. Consider $R=\{x: x\notin x\}$, the set of all sets
that are not elements of themselves, and ask whether $R\in R$. If $R\in R$, then $R$ satisfies the
condition defining $R$, so $R\notin R$. If $R\notin R$, then $R$ satisfies the condition, so $R\in R$.
Either answer contradicts itself. The way out is to restrict which collections count as sets; in
particular, the collection of all sets is not one. Nothing in ECN115 comes close to these limits, which is
why ``set'' can safely be used informally.
\end{external}

\section{Equality of sets}\label{sec:set-equality}

\begin{definition}[Equal sets][def:set-equal]
Two sets $A$ and $B$ are \term{equal}, written $A=B$, if every element of $A$ is also an element of $B$,
and every element of $B$ is an element of $A$. \src{L1, slide 10}
\end{definition}

In the language of \cref{ch:logic}, $A=B$ means: for every object $x$, $x\in A\Leftrightarrow x\in B$. Sets
are equal when they have exactly the same elements, however they are described. For example,
$\{x\in\R: x^2=1\}=\{-1,1\}$, and the set of even prime numbers equals $\{2\}$.

\section{Subsets}\label{sec:subsets}

\begin{definition}[Subset, superset][def:subset]
If each element of $A$ is also an element of $B$, then we write $A\subseteq B$ and call $A$ a
\term{subset} of $B$, and $B$ a \term{superset} of $A$. \src{L1, slide 11}
\end{definition}

With quantifiers: $A\subseteq B$ means $\forall x\in A,\ x\in B$. To prove $A\subseteq B$ you take an
arbitrary element of $A$ and show it lies in $B$. To disprove it you find one element of $A$ that is not
in $B$.

\begin{sourcenote}[Notation: $\subseteq$ and $\subset$]
The slides note that some textbooks write $A\subset B$ for ``$A$ is a subset of $B$''. Be aware that many
other texts use $A\subset B$ to mean a \emph{proper} subset (a subset with $A\ne B$). This book always
writes $\subseteq$ for ``subset'', which is unambiguous. When reading another source, check which
convention it uses.
\end{sourcenote}

\begin{proposition}[Every set is a subset of itself][prop:self-subset]
For every set $A$, $A\subseteq A$. \src{L1, slide 12}
\end{proposition}
\emph{Proof.} Every element of $A$ is an element of $A$. \hfill$\square$

\begin{core}[Equality is two subset relations]
\[
A=B\quad\Longleftrightarrow\quad\big(A\subseteq B\ \text{ and }\ B\subseteq A\big).
\]
To prove that two sets are equal, prove that each is a subset of the other. \src{L1, slide 11}
\end{core}

\begin{proposition}[Equality via subsets][prop:double-inclusion]
For sets $A$ and $B$, $A=B$ if and only if $A\subseteq B$ and $B\subseteq A$.
\end{proposition}
\emph{Proof.} By \cref{def:set-equal}, $A=B$ means ``every element of $A$ is an element of $B$'' \emph{and}
``every element of $B$ is an element of $A$''. By \cref{def:subset}, the first part says exactly
$A\subseteq B$ and the second says exactly $B\subseteq A$. So the two conditions are the same statement
written in different words, and each implies the other. \hfill$\square$

\begin{external}[The empty set is a subset of every set]
For every set $A$, $\emptyset\subseteq A$. To disprove it we would need an element of $\emptyset$ that is not
in $A$, and $\emptyset$ has no elements at all. A statement of the form ``every element of $\emptyset$ has
property $P$'' is therefore true whatever $P$ is; such statements are called \emph{vacuously true}.
\end{external}

\section{Intersection and union}\label{sec:intersection-union}

\begin{definition}[Intersection and union][def:cap-cup]
\begin{itemize}
  \item The collection of objects in both set $A$ and set $B$ is written $A\cap B$ and called the
  \term{intersection} of the two sets.
  \item The collection of objects in either set $A$, or set $B$, or in both sets $A$ and $B$ is written
  $A\cup B$ and called the \term{union} of the two sets.
\end{itemize}
\src{L1, slide 11}
\end{definition}

In symbols, $x\in A\cap B\Leftrightarrow (x\in A\text{ and }x\in B)$, and
$x\in A\cup B\Leftrightarrow(x\in A\text{ or }x\in B)$, where ``or'' is inclusive. Intersection corresponds
to ``and'', union to ``or''. Because ``$P$ and $Q$'' means the same as ``$Q$ and $P$'', and similarly for
``or'', both operations are \emph{commutative}:
\[
B\cap A=A\cap B,\qquad B\cup A=A\cup B. \tag*{\srcin{L1, slide 11}}
\]

\begin{figure}[ht]
\centering
\begin{tikzpicture}[scale=0.82]
  \def\setA{(0,0) circle (1.25)}
  \def\setB{(1.5,0) circle (1.25)}
  % intersection
  \begin{scope}[shift={(0,0)}]
    \begin{scope}\clip \setA; \fill[accent!35] \setB; \end{scope}
    \draw[thick] \setA; \draw[thick] \setB;
    \node at (-0.45,1.55) {$A$}; \node at (1.95,1.55) {$B$};
    \node[font=\small] at (0.75,-1.8) {$A\cap B$};
  \end{scope}
  % union
  \begin{scope}[shift={(4.6,0)}]
    \fill[accent!35] \setA; \fill[accent!35] \setB;
    \draw[thick] \setA; \draw[thick] \setB;
    \node at (-0.45,1.55) {$A$}; \node at (1.95,1.55) {$B$};
    \node[font=\small] at (0.75,-1.8) {$A\cup B$};
  \end{scope}
  % difference
  \begin{scope}[shift={(9.2,0)}]
    \begin{scope}[even odd rule]\clip \setB (-2,-2) rectangle (3.5,2); \fill[accent!35] \setA; \end{scope}
    \draw[thick] \setA; \draw[thick] \setB;
    \node at (-0.45,1.55) {$A$}; \node at (1.95,1.55) {$B$};
    \node[font=\small] at (0.75,-1.8) {$A\setminus B$};
  \end{scope}
\end{tikzpicture}
\caption{Venn diagrams. Each circle represents a set; the shaded region is the set named below it. Left:
the intersection $A\cap B$ (in both). Centre: the union $A\cup B$ (in at least one). Right: the set
difference $A\setminus B$ (in $A$ but not in $B$), defined in \cref{sec:difference}.}
\label{fig:venn}
\end{figure}

\begin{warning}[A picture is not a proof]
Venn diagrams such as \cref{fig:venn} are excellent for seeing what a statement says and for guessing
whether it is true. They are not proofs: a proof must argue from the definitions, for an arbitrary element.
\end{warning}

\begin{example}[Preferences for fruit][ex:fruit]
Emma, Bob, Hannah and Sam like the following fruit:
\begin{align*}
E&=\{\text{oranges},\text{strawberries}\}, & B&=\{\text{oranges},\text{apples},\text{bananas}\},\\
H&=\{\text{bananas},\text{blueberries},\text{apples}\}, & S&=\{\text{bananas},\text{oranges},\text{apples},\text{blueberries}\}.
\end{align*}
Find $B\cap H$, $B\cup H$, $S\cap E$, $B\cap E$ and $E\cap H$, and decide whether $B\subseteq S$ and
$H\subseteq S$. \src{L1, slide 12}
\solution
$B\cap H$ contains the fruit that both Bob and Hannah like: $B\cap H=\{\text{apples},\text{bananas}\}$.

$B\cup H$ contains the fruit liked by at least one of them:
\[B\cup H=\{\text{oranges},\text{apples},\text{bananas},\text{blueberries}\}.\]
This is exactly $S$: as sets
(order irrelevant) the two lists contain the same four fruit, so $B\cup H=S$.

$S\cap E=\{\text{oranges}\}$ and $B\cap E=\{\text{oranges}\}$, so $S\cap E=B\cap E$.

$E\cap H=\emptyset$: Emma and Hannah have no fruit in common.

Every fruit Bob likes is liked by Sam, so $B\subseteq S$; likewise $H\subseteq S$. This also follows from
$B\cup H=S$, since every set is a subset of its union with another set. Finally, every set is a subset of
itself: $E\subseteq E$, $B\subseteq B$ and so on.
\end{example}

\begin{external}[Further algebra of $\cap$ and $\cup$]
Because $\cap$ and $\cup$ mirror ``and'' and ``or'', they obey the same laws:
\begin{itemize}
  \item associativity: $(A\cap B)\cap C=A\cap(B\cap C)$ and $(A\cup B)\cup C=A\cup(B\cup C)$;
  \item distributivity: $A\cap(B\cup C)=(A\cap B)\cup(A\cap C)$ and $A\cup(B\cap C)=(A\cup B)\cap(A\cup C)$;
  \item $A\cap B\subseteq A\subseteq A\cup B$.
\end{itemize}
Each is proved by taking an arbitrary $x$ and translating membership into ``and'' and ``or''.
\end{external}

\section{Set-builder notation}\label{sec:set-builder}

Listing elements works for small sets but not for large or infinite ones. Usually a set is described by
a property that its elements share.

\begin{definition}[Set-builder notation][def:set-builder]
Suppose that a set $A$ might contain some elements that have property $P$. Then
\[
\{x\in A\mid x\text{ has property }P\}
\]
denotes the subset of $A$ that consists of all elements of the set $A$ that have property $P$.
\src{L1, slide 13}
\end{definition}

The vertical bar is read ``such that''. The set $A$ in front of the bar says where the elements are taken
from; the condition after the bar says which of them are kept. The result is always a subset of $A$.

\begin{example}[Reading set-builder notation]
\begin{enumerate}
  \item If $A$ is the set of students at QMUL and $P$ is the property of being born after 2005, then
  $\{x\in A\mid x\text{ was born after 2005}\}$ is the set of students at QMUL who were born after 2005.
  \src{L1, slide 13}
  \item $\{n\in\N\mid n\text{ is even}\}=\{2,4,6,\dots\}$.
  \item $\{x\in\R\mid x^2<4\}$ is the set of real numbers strictly between $-2$ and $2$; in the interval
  notation of \cref{sec:intervals} it is $(-2,2)$.
  \item $\{x\in\R\mid x^2=-1\}=\emptyset$, since no real number squares to $-1$.
\end{enumerate}
\end{example}

\begin{note}[Solving is describing a set]
``Find all $x$ that satisfy $2x+1>0$'' is the same request as ``describe the set
$\{x\in\R\mid 2x+1>0\}$ more simply''. The answer $(-\tfrac12,\infty)$ in \cref{ex:lin-ineq} is a
simpler name for the same set.
\end{note}

\section{Set difference}\label{sec:difference}

\begin{definition}[Set difference][def:difference]
Given two sets $A$ and $B$, the set
\[
A\setminus B=\{x\in A\mid x\notin B\}
\]
is called the \term{difference} between sets $A$ and $B$; we also say that $A\setminus B$ is the result of
\term{subtraction} of the set $B$ from the set $A$. Thus $A\setminus B$ is the subset of elements of $A$
that are not elements of $B$. \src{L1, slide 14}
\end{definition}

\begin{example}[Set differences]
Compute
\begin{enumerate}
  \item $\{\text{apples},\text{oranges},\text{blueberries}\}\setminus\{\text{bananas},\text{oranges}\}$;
  \item $\{2,4,6,8,10,12,14,\dots\}\setminus\{3,6,9,12,15,\dots\}$.
\end{enumerate} \src{L1, slide 14}
\solution
(a) Keep the elements of the first set that are not in the second. Apples and blueberries are not in the
second set; oranges are. So the difference is $\{\text{apples},\text{blueberries}\}$. Bananas play no role:
removing something that was never there changes nothing.

(b) From the even natural numbers remove the multiples of $3$. The even numbers that are multiples of $3$
are the multiples of $6$, so we remove $6,12,18,\dots$ and are left with $\{2,4,8,10,14,16,20,\dots\}$, the
even numbers not divisible by $3$.
\end{example}

\begin{proposition}[Subset and difference][prop:subset-diff]
For sets $A$ and $B$: $A\subseteq B$ implies $A\setminus B=\emptyset$. \src{L1, slide 14} The converse is
also true, so $A\subseteq B\Leftrightarrow A\setminus B=\emptyset$.
\end{proposition}
\emph{Proof.} ($\Rightarrow$) Suppose $A\subseteq B$. An element of $A\setminus B$ would be an element of $A$
that is not in $B$; but every element of $A$ is in $B$. So $A\setminus B$ has no elements:
$A\setminus B=\emptyset$.
($\Leftarrow$) Suppose $A\setminus B=\emptyset$ and take any $x\in A$. If $x\notin B$ we would have
$x\in A\setminus B$, which is impossible; hence $x\in B$. Since $x$ was arbitrary, $A\subseteq B$.
\hfill$\square$

\begin{warning}[Set difference is not commutative]
Unlike $\cap$ and $\cup$, the order matters: in general $A\setminus B\ne B\setminus A$. For example, with
$A=\{1,2\}$ and $B=\{2,3\}$ we have $A\setminus B=\{1\}$ but $B\setminus A=\{3\}$.
\end{warning}

\begin{external}[Two useful facts about $\setminus$]
For any sets $A$ and $B$:
\begin{itemize}
  \item $A\setminus B\subseteq A$, and $A\setminus B$ has no element in common with $B$.
  \item $A\setminus B=A$ if and only if $A\cap B=\emptyset$ (removing $B$ removes nothing exactly when $A$
  and $B$ share no element).
\end{itemize}
\end{external}

\section{Sets in examination questions}\label{sec:sets-exam}

The January 2024 ECN115 paper opened with a sixteen-mark question on sets, consisting of a definition,
two ``does such a thing exist?'' parts and one ``find all'' part (as summarised in the study guide). The
following worked example reproduces that structure.

\begin{example}[Questions on set difference][ex:jan2024]
\begin{enumerate}
  \item Explain the notion of set difference $A\setminus B$.
  \item Are there non-empty sets of natural numbers $A$ and $B$ with $A\setminus B=A$? If yes, give an
  example; if no, argue why.
  \item Are there non-empty sets of natural numbers $A$ and $B$ with $A\setminus B=B$? If yes, give an
  example; if no, argue why.
  \item Find all sets $D\subseteq\N$ such that $\{1,2\}\cup D=\{1,2,3\}$.
\end{enumerate}
\src{SG: Jan 2024 Q1}
\solution
(a) For sets $A$ and $B$, $A\setminus B=\{x\in A\mid x\notin B\}$: the set of elements of $A$ that are not
elements of $B$.

(b) Yes. Take $A=\{1\}$ and $B=\{2\}$. Both are non-empty sets of natural numbers, $1\notin B$, so
$A\setminus B=\{1\}=A$. (Any two non-empty sets with no common element work.)

(c) No. Suppose $A\setminus B=B$ with $B$ non-empty, and take some $b\in B$. Then $b\in A\setminus B$, so by
the definition of set difference $b\notin B$. This contradicts $b\in B$. Hence there are no such sets. (The
argument shows more: $A\setminus B=B$ forces $B=\emptyset$.)

(d) Since $3\in\{1,2,3\}=\{1,2\}\cup D$ and $3\notin\{1,2\}$, we need $3\in D$. Since every element of $D$
lies in $\{1,2\}\cup D=\{1,2,3\}$, we need $D\subseteq\{1,2,3\}$. The elements $1$ and $2$ may or may not be
in $D$, because they are in the union anyway. Hence
\[
D\in\big\{\{3\},\ \{1,3\},\ \{2,3\},\ \{1,2,3\}\big\},
\]
four sets in all, and each of them does satisfy the condition.
\end{example}

\begin{examtip}[``If yes, give an example; if no, argue why'']
This phrasing appears on several past ECN115 papers. A ``yes'' answer needs one explicit example,
checked against every condition (including ``non-empty''). A ``no'' answer needs a general argument,
usually by contradiction: assume such objects exist and derive something impossible. Decide which it is
before writing; a quick experiment with small sets usually tells you.
\end{examtip}

\section{Sets in economics}\label{sec:sets-econ}

\begin{example}[A budget set][ex:budget-set]
A consumer with income $m=100$ buys quantities $x\ge0$ and $y\ge0$ of two goods priced $p_x=5$ and
$p_y=10$. Describe the bundles she can afford as a set, and decide whether the bundles $(10,5)$ and
$(12,5)$ are affordable.
\solution
A bundle is a pair $(x,y)$ of real numbers. The affordable bundles form the \emph{budget set}
\[
B=\{(x,y)\in\R^2\mid x\ge0,\ y\ge0,\ 5x+10y\le100\},
\]
where $\R^2$ denotes the set of all pairs of real numbers. The bundle $(10,5)$ costs $5\cdot10+10\cdot5=100$,
so $(10,5)\in B$ (it is on the budget line). The bundle $(12,5)$ costs $110>100$, so $(12,5)\notin B$.

If a rationing scheme also limits the first good to at most $8$ units, the bundles that are both
affordable and permitted form the intersection $B\cap\{(x,y)\in\R^2\mid x\le 8\}$. Each extra constraint
on a decision-maker corresponds to intersecting with one more set.
\end{example}

% --------------------------------------------------------------------
\section{Chapter review}
% --------------------------------------------------------------------
\subsection*{Summary}
A set is determined by its elements; order and repetition are irrelevant, and the empty set $\emptyset$ has
no elements. $A\subseteq B$ means every element of $A$ is in $B$; every set is a subset of itself; and two
sets are equal exactly when each is a subset of the other. Intersection ($\cap$, ``and''), union ($\cup$,
``or'') and difference ($\setminus$, ``and not'') build new sets; the first two are commutative, the third is
not. Set-builder notation $\{x\in A\mid P\}$ describes the elements of $A$ with property $P$, and solving an
equation or inequality amounts to describing such a set in simpler terms.

\subsection*{Key definitions}
\begin{itemize}
  \item $x\in A$: $x$ is an element of $A$; \ $x\notin A$: it is not.
  \item $A=B$: every element of $A$ is in $B$ and every element of $B$ is in $A$.
  \item $A\subseteq B$: every element of $A$ is in $B$ ($A$ a subset, $B$ a superset).
  \item $A\cap B$: objects in both; \ $A\cup B$: objects in $A$, in $B$ or in both.
  \item $\{x\in A\mid P\}$: elements of $A$ with property $P$.\ \ $A\setminus B=\{x\in A\mid x\notin B\}$.
\end{itemize}

\subsection*{Key facts}
\[
A=B\Leftrightarrow(A\subseteq B\text{ and }B\subseteq A),\quad A\subseteq A,\quad
A\cap B=B\cap A,\quad A\cup B=B\cup A,\quad A\subseteq B\Rightarrow A\setminus B=\emptyset.
\]

\subsection*{Connections}
Membership statements are propositions (\cref{ch:logic}), and $\cap$, $\cup$ correspond to ``and'' and
``or''. The number systems $\N,\Z,\Q,\R$ are sets, one inside the next (\cref{ch:numbers}); intervals are
subsets of $\R$ described in set-builder notation (\cref{sec:intervals}); and the solution of every
inequality in \cref{ch:ineq,ch:abs} is a set, often a union of intervals.

\subsection*{Common mistakes}
\begin{itemize}
  \item Confusing $\in$ (element of) with $\subseteq$ (subset of): $1\in\{1,2\}$ but $\{1\}\subseteq\{1,2\}$.
  \item Assuming $A\setminus B=B\setminus A$.
  \item Forgetting that ``or'' in $A\cup B$ includes elements in both sets.
  \item Treating a Venn diagram as a proof.
  \item In ``find all sets'' questions, listing only one set when several satisfy the conditions.
\end{itemize}

\subsection*{Exam checklist}
\begin{checklist}
  \item Define subset and set equality, and state $A=B\Leftrightarrow(A\subseteq B$ and $B\subseteq A)$.
  \item Define set difference in set-builder notation.
  \item Compute intersections, unions and differences of given sets.
  \item Answer ``does such a set exist?'' questions with an example or an argument by contradiction.
  \item Find all sets satisfying given union, intersection and subset conditions.
\end{checklist}

% --------------------------------------------------------------------
\section{Practice questions}
% --------------------------------------------------------------------
\pq[basic] Let $A=\{1,2,3,4,5\}$ and $B=\{4,5,6\}$. Find $A\cap B$, $A\cup B$, $A\setminus B$ and $B\setminus A$.

\pq[basic] List the elements of (a) $\{n\in\N\mid n<5\}$; (b) $\{n\in\N\mid n^2\le 30\}$;
(c) $\{x\in\R\mid x^2=4\}$; (d) $\{x\in\R\mid x^2+1=0\}$.

\pq[conceptual] Let $C=\{\{1,2\},\{1,3\},\{2,3\}\}$. Which of the following are true?
(a) $\{1,2\}\in C$; (b) $1\in C$; (c) $\{2,1\}\in C$; (d) $\{\{1,2\}\}\subseteq C$; (e) $\{1,2\}\subseteq C$.

\pq[mathematical] Prove that for any sets $A$ and $B$, $A\cap B\subseteq A$ and $A\subseteq A\cup B$.

\pq[mathematical] Prove that if $A\subseteq B$ and $B\subseteq C$, then $A\subseteq C$.

\pq[application] A firm can produce any whole number of units from $0$ to $50$ inclusive. A regulator
requires output of at least $20$ units, and capacity constraints in a second plant exclude the outputs
$30$ to $35$ inclusive. Describe the set of permitted outputs using set-builder notation and set operations,
and state how many outputs are permitted.

\pq[multi-step] Are there non-empty sets of natural numbers $A$ and $B$ such that $A\cup B=A\cap B$? If yes,
give an example and describe all such pairs; if no, argue why.

\pq[exam-style] (a) Explain what it means for a set $A$ to be a subset of a set $B$, and state what it means
for two sets to be equal in terms of subsets.
(b) Find all sets $D$ of natural numbers that satisfy both $\{2,4\}\cap D=\{2\}$ and $D\subseteq\{1,2,3,4\}$.
(c) Let $A=\{x\in\N\mid x\le10\}$ and $B=\{x\in\N\mid x\text{ is even}\}$. List the elements of $A\setminus B$
and decide whether $A\setminus B=B\setminus A$.

% --------------------------------------------------------------------
\section{Solutions to practice questions}
% --------------------------------------------------------------------
\pqsol{2.1}
$A\cap B=\{4,5\}$; $A\cup B=\{1,2,3,4,5,6\}$; $A\setminus B=\{1,2,3\}$; $B\setminus A=\{6\}$. Note that
$A\setminus B\ne B\setminus A$.

\pqsol{2.2}
(a) $\{1,2,3,4\}$ (recall $0\notin\N$). (b) $n^2\le30$ holds for $n=1,2,3,4,5$ ($5^2=25$) but not $n=6$
($36$); so $\{1,2,3,4,5\}$. (c) $\{-2,2\}$. (d) $\emptyset$.

\pqsol{2.3}
(a) True. (b) False: the elements of $C$ are sets, not numbers. (c) True: $\{2,1\}=\{1,2\}$, since order does
not matter. (d) True: the only element of $\{\{1,2\}\}$ is $\{1,2\}$, which is in $C$. (e) False: $1$ is an
element of $\{1,2\}$ but $1\notin C$.

\pqsol{2.4}
Let $x\in A\cap B$ be arbitrary. Then $x\in A$ and $x\in B$; in particular $x\in A$. So $A\cap B\subseteq A$.
Now let $x\in A$ be arbitrary. Then ``$x\in A$ or $x\in B$'' is true, so $x\in A\cup B$. So $A\subseteq A\cup B$.

\pqsol{2.5}
Let $x\in A$ be arbitrary. Since $A\subseteq B$, $x\in B$. Since $B\subseteq C$, $x\in C$. As $x$ was an
arbitrary element of $A$, $A\subseteq C$.

\pqsol{2.6}
Let $X=\{n\in\Z\mid 0\le n\le 50\}$ be the feasible outputs. The permitted outputs are
\[
\{n\in X\mid n\ge 20\}\setminus\{n\in\Z\mid 30\le n\le 35\}=\{20,21,\dots,29\}\cup\{36,37,\dots,50\}.
\]
There are $10+15=25$ permitted outputs.

\pqsol{2.7}
Yes: $A=B=\{1\}$ gives $A\cup B=\{1\}=A\cap B$. In fact this happens exactly when $A=B$. Always
$A\cap B\subseteq A\subseteq A\cup B$ (Question~2.4). If $A\cup B=A\cap B$, the outer two sets coincide, so
$A$ is squeezed between equal sets and $A=A\cap B=A\cup B$; the same argument gives $B=A\cup B$; hence
$A=B$. Conversely, if $A=B$ then $A\cup B=A=A\cap B$.

\pqsol{2.8}
(a) $A\subseteq B$ means that every element of $A$ is also an element of $B$. Two sets are equal,
$A=B$, exactly when $A\subseteq B$ and $B\subseteq A$.
(b) $\{2,4\}\cap D=\{2\}$ says $2\in D$ and $4\notin D$. $D\subseteq\{1,2,3,4\}$ says $D$ has no other elements.
So $1$ and $3$ are each free to be in or out: $D\in\{\{2\},\{1,2\},\{2,3\},\{1,2,3\}\}$.
(c) $A=\{1,\dots,10\}$, so $A\setminus B=\{1,3,5,7,9\}$. But $B\setminus A=\{12,14,16,\dots\}$, the even numbers
greater than $10$. For example $12\in B\setminus A$ while $12\notin A\setminus B$, so $A\setminus B\ne B\setminus A$.
